\documentclass[openright, oneside, a4paper, article, brazil]{abntex2}
\usepackage{formularios}
% --------- %
% Metadados %
% --------- %
\hypersetup{%
    colorlinks=false,
    allbordercolors=white,
    pdftitle = {TÍTULO DO FORMULÁRIO},
    pdfauthor = {Conselho Regional de Psicologia de São Paulo},
    pdfsubject = {ASSUNTO DO FORMULÁRIO},
    pdfkeywords = {%
            palavra-chave1,
            palavra-chave2,
            palavra-chave3,
        }%
}
\title{TÍTULO DO FORMULÁRIO}
\author{Conselho Regional de Psicologia de São Paulo}
% --------- %
% DOCUMENTO %
% --------- %
\begin{document}
\pagestyle{crp}
\centering
\Form
\section*{TÍTULO DO FORMULÁRIO}%
% -------------------- %
% EXEMPLOS DE CONTEÚDO %
% -------------------- %
% -------------------------- %
% Identificação/qualificação %
% -------------------------- %
\setasuspacing{\DoubleSpacing}
% Define espaçamento duplo para que as caixas de texto não fiquem sobrepostas
\justifying
% Parágrafo justificado
\preenctext{n_cnpj}{15em}{Nº CNPJ }
% Caixa de texto %
%  Sintaxe: \preenctext{#1}{#2}{#3}
%   #1: id único para caixa de texto;
%   #2: extensão da caixa de texto, com unidade de medida (p. ex., "em", "cm"ou "pt");
%   #3: texto que precede a caixa (incluir espaço após o texto para que não fique colado na caixa).

\preenctext{denominacao}{38.5em}{Denominação social \hfill}\\%
\preenctext{denominacao2}{\linewidth}{}\\%
% \\: fim de linha (*soft break*)
% \hfill: preenche espaço horizontal entre o texto e a caixa
% \linewidth: unidade de medida equivalente à largura da linha que está sendo utilizada;
% diferente de \textwidth, que é a unidade equivalente à largura da linha na mancha gráfica
\vfill
% Preenche espaço vertical

\setasuspacing{\SingleSpacing} % Define espaçamento simples
\begin{framed}
  % Cria moldura em todos os ambientes que estiverem entre o ambiente framed (i.e., entre \begin{framed} e \end{framed}).
  \centering
  % Centraliza o parágrafo
  \vspace*{-3ex}
  % Reduz espaçamento vertical em 3 ex (altura da letra "x") antes de \subsection
  \subsection{Falta}
  \noindent
  % Sem recuo
  \small
  % Tamanho da fonte
  Faltou no dia \dataabrev{12}{1} ou no período de \dataabrev{12}{2} a \dataabrev{12}{3}
  % Data abreviada %
  %    Sintaxe:
  %      \dataabrev{#1}{#2}
  %        #1: tamanho da fonte em pontos (pt);
  %        #2: número de id único para cada data.
\end{framed}
\begin{framed}%
  \centering\vspace*{-3ex}
  \subsection{Falta de marcação}
  \noindent\small Não marcou dia \dataabrev{12}{5}
  \vspace{\baselineskip}% Espaço vertical de uma linha-base

  \footnotesize % Tamanho da fonte
  \begin{minipage}{0.25\linewidth}
    % Começa uma minipágina com 1/4 da extensão da linha
    \preenctext{12}{entrada-manha}{6em}{Entrada manhã: }
  \end{minipage}%
  % Uso de "%" após o fim de minipágina para que a próxima minipágina fique na mesma linha, não abaixo
  \begin{minipage}{0.25\linewidth}
    \preenctext{12}{saida-almoco}{6em}{Saída almoço: }
  \end{minipage}%
  \begin{minipage}{0.25\linewidth}
    \preenctext{12}{retorno-almoco}{6em}{Retorno almoço: }
  \end{minipage}%
  \begin{minipage}{0.25\linewidth}
    \preenctext{12}{saida-tarde}{6em}{Saída tarde: }
  \end{minipage}
\end{framed}
\vspace{\baselineskip}

\end{document}
